\documentclass[../main.tex]{subfiles}
\graphicspath{{\subfix{../images/}}}
\begin{document}

\subsection{Curvilinear Coordinates} \label{sec:curvy_coords}
Extending the Continuum Mechanics formulation to the formulation of shells,
 it becomes impossible to describe them using an orthogonal basis system for obvious reasons.
 Thus it is necessary to use a curvilinear coordinate system to account for the bent surface.
 Although shell model are often formulated using local cartesian coordinate systems as the
 classical relations described in the previous section can be extended to account for shells.
 However, the formulation under investigation feels unnecessary to make this detour and the relevant
 mathematics is described using curvilinear coordinates.

The membrane surface, denoted $\mathcal{S}$, is fully characterised by 
 the parametric description,
 \begin{equation} \label{parametric_definition}
    \boldsymbol{x} = \boldsymbol{x}(\xi^{1}, \xi^{2}) \, .
 \end{equation}
This is the mappinng of the point $(\xi^{1}, \xi^{2})$ in the parameter domain $\mathcal{P}$
 to the material point $\boldsymbol{x} \in \mathcal{S}$.
According to the non-standard Tensor notation used in Continuum Mechanics,
 Greek indices take on the values 1 and 2 on the surface. And as in the previous section,
 lower case Latin indices take the values 1, 2, and 3, in Euclidean space and summation is
 implied on repeated indices.
The tangent vectors that form the basis at a coordinate $\xi^(\alpha)$ at point $\boldsymbol{x} \in \mathcal{S}$ are
 given by,
 \begin{equation}\label{tangent_parametric}
    \boldsymbol{a}_{\alpha} = \frac{\partial \boldsymbol{x}}{\partial \xi^{\alpha}}, \,\,\, \alpha = 1, 2
 \end{equation}
In general, they are not orthonomal. This basis is thus characterised by the covariant metric tensor,
\begin{equation}\label{metric_def}
    a_{\alpha \beta} := \boldsymbol{a}_{\alpha} \cdot \boldsymbol{a}_{\beta}.
\end{equation}
In case of an orthonormal basis system, the covariant metric tensor reduces to the Identity tensor.

The covariant components of the metric tensor can be obtained by its inversion,
\begin{equation}\label{contravar_metric_def}
    [a^{\alpha \beta}] := [a_{\alpha \beta}]^{-1}.
\end{equation}
This enables the construction of a dual basis - the contra-variant counterparts of the covariant base vectors, $\boldsymbol{a}_{\alpha}$,
\begin{align}\label{contravar_def}
    \boldsymbol{a}^{\alpha} &:= a^{\alpha \beta} \boldsymbol{a}_{\beta},\\
    \boldsymbol{a}^{\alpha} &\cdot \boldsymbol{a}^{\beta} = a^{\alpha \beta},\\
    \boldsymbol{a}_{\alpha} &\cdot \boldsymbol{a}^{\beta} = \delta^{\beta}_{\alpha},
\end{align}
where $\delta^{\beta}_{\alpha}$ is the Kronecker delta.
%
The unit normal at a point $\boldsymbol{x} \in \mathcal{S}$ is defined by,
\begin{equation}\label{normal}
    \boldsymbol{n} = \frac{\boldsymbol{a}_{1} \times \boldsymbol{a}_{2}}{\| \boldsymbol{a}_{1} \times \boldsymbol{a}_{2} \|}
    = \frac{\boldsymbol{a}_{1} \times \boldsymbol{a}_{2}}{\sqrt{\det a_{\alpha \beta}}}.
\end{equation}

Any vector, $\boldsymbol{v}$ on $\mathcal{S}$ can thus be decomposed using the covariant or the contravariant basis,
\begin{equation}
    \boldsymbol{v} = v^{\alpha} \boldsymbol{a}_{\alpha} + v_{\mathrm{n}}\boldsymbol{n} = v_{\alpha} \boldsymbol{a}^{\alpha} + v_{\mathrm{n}},
\end{equation}
 where $v_{\alpha}$ and $v^{\alpha}$ denote the covariant and the contravariant components of the vector respectively.
The derivative and the co-variant derivative of the tangent vectors, $\boldsymbol{a}_{\alpha}$, can be defined as,
\begin{align}\label{tangent_derivatives}
    \boldsymbol{a}_{\alpha , \beta} &= \frac{\partial \boldsymbol{a}_{\alpha}}{\partial \xi^{\beta}},\\
    \boldsymbol{a}_{\alpha ; \beta} := \boldsymbol{a}_{\alpha , \beta} - &(\boldsymbol{a}_{\alpha , \beta} \cdot \boldsymbol{a}^{\gamma}) \boldsymbol{a}_{\gamma} = \boldsymbol{a}_{\alpha , \beta} - \Gamma^{\gamma}_{\alpha \beta} \boldsymbol{a}_{\gamma},
\end{align}
where $\Gamma^{\gamma}_{\alpha \beta}$ are the Christoffel symbols of the second kind.
Defining the identity tensor on the surface, $\mathcal{S}$,
\begin{equation}\label{surface_identity}
    \boldsymbol{1} = \boldsymbol{a}_{\alpha} \otimes \boldsymbol{a}^{\alpha} = \boldsymbol{a}^{\alpha} \otimes \boldsymbol{a}_{\alpha} = \boldsymbol{\tilde{1}} - \boldsymbol{n} \otimes \boldsymbol{n}, 
\end{equation}
where $\boldsymbol{\tilde{1}}$ is the usual identity tensor in $\mathbb{R}^{3}$, two results can be derived.
\begin{equation}\label{con_tgnt}
    \boldsymbol{a}_{\alpha ; \beta} = (\boldsymbol{n} \otimes \boldsymbol{n}) \boldsymbol{a}_{\alpha,\beta}.
\end{equation}
And the definition of the curvature tensor with covariant components \footnote{The proofs to \eqref{con_tgnt} and \eqref{curvature_tensor_def} can be found in the appendix, \eqref{appendices:cov_tangent_derivative} and \eqref{appendices:curvature_tensor_derivation}.}
 , $\boldsymbol{b} = b_{\alpha \beta} \boldsymbol{a}^{\alpha}
 \otimes \boldsymbol{a}^{\beta}$,
\begin{equation} \label{curvature_tensor_def}
    b_{\alpha \beta} := \boldsymbol{n} \cdot \boldsymbol{a}_{\alpha ; \beta} = \boldsymbol{n} \cdot \boldsymbol{a}_{\alpha,\beta}.
\end{equation}
\subsection{Kinematics}
%---------------------------------------------------------------------------------
As defined in Section \ref{sec:continuum_kinematics}, to consider the deformation of the membrane surface, we distinguish between the reference configuration,
 $\mathcal{S}_{0}$, and the current configuration, $\mathcal{S}$. To describe quantities on $\mathcal{S}$,
 are written in lowercase symbols - $\boldsymbol{x}, \boldsymbol{a}_{\alpha}, \boldsymbol{a}^{\alpha}, \boldsymbol{n}, a_{\alpha \beta}, b_{\alpha \beta}$.
 Correspondingly, uppercase symbols are used to describe the quantities on $\mathcal{S}_{0}$.
 The deformation of a line element in $\mathcal{S}_{0}$ can be described using the Deformation tensor,
\begin{equation}\label{curvy_F}
    \begin{split}
    \mathrm{d}\boldsymbol{x} &= {\boldsymbol{a}_{\alpha} \otimes \boldsymbol{A}^{\alpha}} \mathrm{d} \boldsymbol{X},\\
    &\text{where } \mathrm{d} \boldsymbol{X} = \frac{\partial \boldsymbol{X}}{\partial \xi^{\alpha}} = \boldsymbol{A}_{\alpha} \mathrm{d} \xi^{\alpha}, \notag\\
    &\text{and } \mathrm{d} \boldsymbol{x} = \boldsymbol{a}_{\alpha} \mathrm{d} \xi^{\alpha}. \notag
\end{split}
\end{equation}

The following transformations can be described using \textbf{F},
\begin{align}\label{a_and_A}
    \boldsymbol{a}_{\alpha} &= \left(\boldsymbol{a}_{\alpha} \otimes \boldsymbol{A}^{\alpha}\right) \boldsymbol{A}_{\alpha} = \boldsymbol{FA}_{\alpha},  &  \boldsymbol{a}^{\alpha} &= \left(\boldsymbol{a}^{\alpha} \otimes \boldsymbol{A}_{\alpha}\right) \boldsymbol{A}^{\alpha} = \boldsymbol{F}^{-T} \boldsymbol{A}^{\alpha},\\
    \boldsymbol{A}_{\alpha} &= \left(\boldsymbol{A}_{\alpha} \otimes \boldsymbol{a}^{\alpha}\right) \boldsymbol{a}_{\alpha} = \boldsymbol{F}^{-1} \boldsymbol{a}_{\alpha},  &   \boldsymbol{A}^{\alpha} &= \left(\boldsymbol{A}^{\alpha} \otimes \boldsymbol{a}_{\alpha}\right) \boldsymbol{a}^{\alpha} = \boldsymbol{F}^{T} \boldsymbol{a}^{\alpha}.
\end{align}

The right and left Cauchy-Green surface tensors are,
\begin{align}\label{curvy_strain}
    \boldsymbol{C} &= \boldsymbol{F}^{T} \boldsymbol{F} = a_{\alpha \beta} \boldsymbol{A}^{\alpha} \otimes \boldsymbol{A}^{\beta},   &   \boldsymbol{C}^{-1} &= (\boldsymbol{F}^{T} \boldsymbol{F})^{-1} = a^{\alpha \beta} \boldsymbol{A}_{\alpha} \otimes \boldsymbol{A}_{\beta},\\
    \boldsymbol{b} &= \boldsymbol{F} \boldsymbol{F}^{T} = A^{\alpha \beta} \boldsymbol{a}_{\alpha} \otimes \boldsymbol{a}_{\beta},  &   \boldsymbol{b}^{-1} &= (\boldsymbol{F} \boldsymbol{F}^{T})^{-1} = A_{\alpha \beta} \boldsymbol{a}^{\alpha} \otimes \boldsymbol{a}^{\beta}.
\end{align}

Similar to Section \ref{sec:continuum_kinematics}, the area elements on the surfaces can be related to each other
 with the help of the parametric coordinates $(\xi^{1}, \xi^{2})$. The differential surface area in the current and reference configuration and their relation can be given as,
 \begin{gather} \label{area_something_man}
    \mathrm{d}a := \Vert (\boldsymbol{a}_{1} \mathrm{d}\xi^{1}) \times ((\boldsymbol{a}_{2} \mathrm{d}\xi^{2})) \Vert = \Vert \boldsymbol{a}_{1} \times \boldsymbol{a}_{2} \Vert\, \mathrm{d}\square = \sqrt{\det a_{\alpha \beta}}\,\, \mathrm{d}\square =: J_{a} \mathrm{d}\square,\\
    \mathrm{d}A := \Vert \boldsymbol{A}_{1} \times \boldsymbol{A}_{2} \Vert \, \mathrm{d}\square = \sqrt{\det A_{\alpha \beta}}\,\, \mathrm{d}\square =: J_{A} \mathrm{d}\square,\\    
    \implies \mathrm{d}a = \frac{J_{a}}{J_{A}} \mathrm{d}A =: J \mathrm{d}A
\end{gather}
 where $\mathrm{d} \square := \mathrm{d} \xi^{1} \, \mathrm{d} \xi^{2}$. 
\subsection{Momentum Balance}
%---------------------------------------------------------------------------------
At any given point, the Cauchy stress tensor $\boldsymbol{\sigma}$ can be written as,
\begin{equation}\label{curvy_stress}
    \boldsymbol{\sigma} = \sigma^{\alpha \beta} \boldsymbol{a}_{\alpha} \otimes \boldsymbol{a}_{\beta} + \sigma^{3\alpha}(\boldsymbol{n}
    \otimes \boldsymbol{a}_{\alpha} + \boldsymbol{a}_{\alpha} \otimes \boldsymbol{n}) + \sigma^{33} \boldsymbol{n} \otimes \boldsymbol{n}.
\end{equation}
As can be seen from \eqref{curvy_stress}, $\boldsymbol{\sigma}$ is symmetric in this formulation. For membranes, it can be assumed that $\sigma^{3\alpha} = \sigma^{33} = 0$,
 which reduces the stress tensor to the \emph{in-plane Cauchy stress tensor}
\begin{equation}\label{curvy_planar_stress}
    \boldsymbol{\sigma} = \sigma^{\alpha \beta} \boldsymbol{a}_{\alpha} \otimes \boldsymbol{a}_{\beta}.
\end{equation}
From the in-plane Cauchy tensor, the internal traction acting on the surface perpendicular to $\boldsymbol{a}^{\alpha}$ can be written as\footnote{The proof to \eqref{traction_def} can be found in the appendix},
\begin{equation}\label{traction_def}
    \boldsymbol{t}^{\alpha} = \boldsymbol{\sigma} \boldsymbol{a}^{\alpha} = \sigma^{\beta \alpha} \boldsymbol{a}_{\beta},
\end{equation}

Given a distributed surface force, \textbf{f},
\begin{equation}\label{surface_force}
    \boldsymbol{f} = f_{\alpha} \boldsymbol{a}^{\alpha} + p \boldsymbol{n} = f^{\alpha} \boldsymbol{a}_{\alpha} + p \boldsymbol{n}.
\end{equation}
the balance of linear momentum in its strong form is (STEIGMANN REFERENCE),
\begin{equation}\label{lin_momentum_curvy}
    \boldsymbol{t}^{\alpha}_{;\alpha} + \boldsymbol{f} = \boldsymbol{0}.\
\end{equation}
Using \eqref{traction_def} and \eqref{curvature_tensor_def}, it can be derived that the covariant derivative of $\boldsymbol{t}^{\alpha}$ can be written as
\begin{equation} \label{traction_cov_derivative}
    \boldsymbol{t}^{\alpha}_{;\alpha} = \sigma ^{\beta \alpha}_{;\alpha}\boldsymbol{a}_{\beta} + \sigma^{\beta \alpha} b_{\beta \alpha} \boldsymbol{n}.
\end{equation}
This enables the decomposition of \eqref{lin_momentum_curvy} into an equation for in-plane and out-of-plane equilibrium.
\begin{equation}\label{lin_momentum_curvy_decomposition}
    \begin{cases}
        \sigma^{\beta \alpha}_{;\alpha} + f^{\beta} = 0 \,\,\, \text{(in-plane)}\\
        \sigma^{\beta \alpha} b_{\beta \alpha} + p = 0 \, \, \, \text{(out-of-plane)}
    \end{cases}
\end{equation}
\linebreak
\subsection{Weak Form} \label{sec:curvy_weak_form}
With the help of a mathematically admissible test function $\boldsymbol{w}$ on $\mathcal{S}$, to derive the weak form \eqref{lin_momentum_curvy} can be written as  
\begin{equation}\label{weak_form_1}
    \int_{\mathcal{S}} \boldsymbol{w} \left[ \boldsymbol{t}_{; \alpha}^{\alpha} + \boldsymbol{f} \right] \, \mathrm{d} a = 0.
\end{equation}
Decomposing $\boldsymbol{w}$ into its surface and normal components     
\begin{equation}\label{weak_form_2}
    0 = \int_{\mathcal{S}} \left[ w_{\gamma} \boldsymbol{a}^{\gamma} + w \boldsymbol{n} \right] \left[ \sigma_{~~; \alpha}^{\beta \alpha} \boldsymbol{a}_{\beta} + \sigma^{\beta \alpha} b_{\beta \alpha} \boldsymbol{n} + \boldsymbol{f} \right] \, \mathrm{d} a
\end{equation}
and using $ \boldsymbol{t} = m_{\alpha} \boldsymbol{t}^{\alpha} = m_{\alpha} \sigma^{\beta \alpha} \boldsymbol{a}_{\beta} $,
 the weak form can be written as \footnote{The complete derivation can be found in the Appendix \ref{appendices:weak_form_derivation}},
    
    \begin{equation}\label{weak_form_curvy}
    0 = \int_{\partial \mathcal{S}} w_{\gamma} \overline{t}^{\gamma} \, \mathrm{d} s - \int_{\mathcal{S}} w_{\gamma ; \alpha} \sigma^{\gamma \alpha} \, \mathrm{d} a + \int_{\mathcal{S}} w_{\gamma} f^{\gamma} \, \mathrm{d} a + \int_{\mathcal{S}} w \sigma^{\beta \alpha} b_{\beta \alpha} \, \mathrm{d} a + \int_{\mathcal{S}} w p \, \mathrm{d} s.
    \end{equation}
    
    And thus it is possible to write
    
    \begin{equation}\label{eq:weak_8}
        G_{int} - G_{ext} = 0 \quad \forall \boldsymbol{w} \in \mathcal{W}
    \end{equation}
    
    with 
    
    \begin{equation}\label{curvy_Gint}
        G_{\mathrm{int}} := \int_{\mathcal{S}} w_{\alpha ; \beta} \sigma^{\alpha \beta} \, \mathrm{d} a - \int_{\mathcal{S}} w \sigma^{\alpha \beta} b_{\alpha \beta} \, \mathrm{d} a 
    \end{equation}
    
    \begin{equation}\label{curvy_Gext}
        G_{\mathrm{ext}} := \int_{\mathcal{S}} w_{\alpha} f^{\alpha} \, \mathrm{d} a + \int_{\partial_{t} \mathcal{S}} w_{\alpha} \overline{t}^{\alpha} \, \mathrm{d} s + \int_{\mathcal{S}} w p \, \mathrm{d} a
    \end{equation}
    
    
    
    With~\eqref{appendices:w_ab} it is possible to rewrite~\eqref{curvy_Gint} as
    
    \begin{equation}\label{eq:weak_11}
    G_{\mathrm{int}} = \int_{\mathcal{S}} \boldsymbol{w}_{; \alpha} \cdot \sigma^{\alpha \beta} \boldsymbol{a}_{\beta} \, \mathrm{d} a
    \end{equation}
%---------------------------------------------------------------------------------
%
%
\subsection{Material Constitution}
For a Neo-Hookean material with a given strain energy density function
\begin{equation}\label{eq:NH_strain}
    W = \frac{\Lambda}{4} \left( J^{2} - 1 - 2 \ln{J} \right) + \frac{\mu}{2} \left( I_{1} - 2 - 2 \ln{J} \right) 
\end{equation}
%
and the traces of the right and left Cauchy-Green tensors in~\eqref{curvy_strain}
\begin{equation}\label{eq:CGtrace}
    I_{1} := \boldsymbol{C} : \boldsymbol{I}_{0} = \boldsymbol{B} : \boldsymbol{I} = a_{\alpha \beta} \boldsymbol{A}^{\alpha} \otimes \boldsymbol{A}^{\beta} : \boldsymbol{A}_{\gamma} \otimes \boldsymbol{A}^{\gamma} = a_{\alpha \beta} \delta^{\alpha}_{\gamma} A^{\beta \delta} = a_{\gamma \beta} A^{\beta \gamma} = a_{\alpha \beta} A^{\alpha \beta}
\end{equation}
the Kirchhoff stress tensor can be written as 
\begin{equation}\label{kirchhoff_stress_def}
        \tau^{\alpha \beta} = 2 \frac{\partial W}{\partial a_{\alpha \beta}} = \left( \frac{\Lambda}{2} \left( J^{2} - 1 \right) - \mu \right) a^{\alpha \beta} - \mu A^{\alpha \beta}.
\end{equation}
Using~\eqref{kirchhoff_stress_def}, the compliance matrix in the current configuration
\begin{equation}\label{curvy_compliance_current_def}
    c^{\alpha \beta \gamma \delta} = \frac{\partial^{2} W}{\partial a_{\alpha \beta} \partial a_{\gamma \delta}} = 2 \frac{\partial \tau^{\alpha \beta}}{\partial a_{\gamma \delta}}
\end{equation}
%
can be written as\footnote{A complete derivation can be found in the Appendix \ref{appendices:compliance_matrix_curvy}}
\begin{equation}\label{curvy_compliance_current}
        c^{\alpha \beta \gamma \delta} = \Lambda J^{2} a^{\alpha \beta} a^{\gamma \delta} + \left( \frac{\Lambda}{2} \left( J^{2} - 1 \right) - \mu \right) m^{\alpha \beta \gamma \delta}.
\end{equation}
where
\begin{equation} \label{mabgd}
    m^{\alpha \beta \gamma \delta} = \frac{1}{a} \left( e^{\alpha \gamma} e^{\delta \beta} + e^{\alpha \delta} e^{\gamma \beta} \right) - 2 a^{\alpha \beta} a^{\gamma \delta}
\end{equation}
%---------------------------------------------------------------------------------
%
%
\subsection{Volume Constraint}
%---------------------------------------------------------------------------------
%
%
\subsection{Discretization}
Following the procedures and nomenclature defined in Section \ref{sec:nonlinear_discretisation}, for quadrilateral finite elements $\Omega^{e}_{0}$ and $\Omega^{e}$
 related to a reference element in the parametric space, the discretisations of the nodal variables in the current and the reference configurations are given according to \eqref{nodal_discretisations}.
The tangent vectors~\eqref{tangent_parametric} can be approximated by
\begin{equation} \label{tangent_parametric_discrete}
    \boldsymbol{a}_{\alpha} \approx \sum^{n_{en}}_{I} N_{I,\alpha} \boldsymbol{x}_{I} = \mathbf{N}_{,\alpha} \mathbf{x}_{e}
\end{equation}
where summation is carried out over the $n_{en}$ nodes of the element. The shape function matrix and its derivative with respect to the parametric coordinates
 is given as
\begin{align} \label{curvy_shape_function}
    N &:= \left[N_{1}\boldsymbol{\tilde{1}}, N_{2} \boldsymbol{\tilde{1}},\dots, N_{n_{en}} \boldsymbol{\tilde{1}}\right],\\
    N_{I,\alpha} & = \frac{\partial N_{I}}{\partial \xi^{\alpha}}.
\end{align}
%
Using~\eqref{eq:w_cov}, and~\eqref{tangent_parametric_discrete}, the internal virtual work
 over a finite element is discretised as
 \begin{equation} \label{curvy_Gint_discrete}
    \begin{split}
    G^{e}_{\text{int}} = \int_{\Omega_{0}^{e}} \boldsymbol{w}_{;\alpha} \cdot \tau^{\alpha \beta} \boldsymbol{a}_{\beta} \,dA 
    &\approx \mathbf{w}^{T}_{e} \int_{\Omega_{0}^{e}} \mathbf{N}^T_{,\alpha} \tau^{\alpha \beta} \mathbf{N}_{,\beta}\,dA \,\, \mathbf{x}_{e}\\ & =: \mathbf{w}^{T}_{e} \mathbf{f}^{e}_{int}.
    \end{split}
 \end{equation}
The internal virtual work~\eqref{curvy_Gint} can further be decomposed and discretised into its in-plane and out-of-plane components
 using~\eqref{eq:w_alp_cov} and the discretisation of the curvature tensor <<MENTION??>> in order to obtain the corresponding force vectors, $\mathbf{f}^{e}_{\text{inti}}$
 and $\mathbf{f}^{e}_{\text{into}}$
\begin{align} \label{curvy_Gint_decomposed_discretised}
    G_{\text{inti}}^{e} &= \int_{\Omega_{0}} w_{\alpha;\beta} \tau^{\alpha \beta}\,dA \approx \mathbf{w}_{e}^{T} \mathbf{f}^{e}_{\text{inti}} = \mathbf{w}_{e}^{T} \int_{\Omega_{0}^{e}} \tau^{\alpha \beta}
    \left[ \mathbf{N}_{, \alpha}^{T} \mathbf{N}_{, \beta} + \mathbf{N}^{T} \left( \boldsymbol{n} \otimes \boldsymbol{n} \right) \mathbf{N}_{, \alpha \beta} \right]\,dA \mathbf{x}_{e},\\
    G_{\text{into}}^{e} &= -\int_{\Omega_{0}^{e}} w \tau^{\alpha \beta} b_{\alpha \beta}\,dA 
    \approx \mathbf{w}_{e}^{T} \mathbf{f}^{e}_{\text{into}} = - \mathbf{w}_{e}^{T}\int_{\Omega_{0}^{e}} \tau^{\alpha \beta}\mathbf{N}^{T} \left( \boldsymbol{n} \otimes \boldsymbol{n} \right) \mathbf{N}_{, \alpha \beta}\,dA \mathbf{x}_{e}.
\end{align}

The external virtual work~\eqref{curvy_Gext} can be discretised to yield the external force vector $\mathbf{f}^{e}_{\text{ext}}$, 
\begin{equation}\label{curvy_Gext_discretised}
    \begin{split}
        G_{\text{ext}}^{e} &= \int_{\Omega_{0}^{e}} \boldsymbol{w} \cdot \boldsymbol{f}_{0}\,dA + \int_{\partial_{t}\Omega^{e}} \boldsymbol{w}\cdot \overline{\boldsymbol{t}}\,ds + \int_{\Omega^{e}} wp \,da,\\
        &\approx \mathbf{w}^{T}_{e} \left(\int_{\Omega_{0}^{e}} \mathbf{N}^{t} \boldsymbol{f}_{0} \,dA + \int_{\partial_{t} \Omega^{e}} \mathbf{N}^{T} \overline{\boldsymbol{t}}\,ds + \int_{\Omega^{e}} \mathbf{N}^{T} p \boldsymbol{n} \,da \right) =: \mathbf{w}^{T}_{e} \mathbf{f}^{e}_{\text{ext}}.
    \end{split}
\end{equation}
In~\eqref{curvy_Gext_discretised}, the external loading is assumed to be of the form $\boldsymbol{f} = \boldsymbol{f}_{0}/J + p\boldsymbol{n}$, where $\boldsymbol{f}_{0}$ and $p$ are the specified loading parameters. 
Using \eqref{curvy_Gint_discrete} and \eqref{curvy_Gext_discretised}, the discretised weak form can be written as
\begin{equation} \label{curvy_weak_form_discretised}
    \mathbf{w}^{T}\left[\mathbf{f}^{e}_{\text{int}} - \mathbf{f}^{e}_{\text{ext}}\right] = \boldsymbol{0}.
\end{equation}
Since the discretised test functions $\mathbf{w}$ are arbitrary and zero on the Dirichlet boundary, the system of discretised nonlinear equilibrium equations
 to be solved for the unknown nodal positions $\mathbf{x}$ is
 \begin{equation} \label{curvy_f_system}
    \mathbf{f} := \mathbf{f}_{\text{int}} - \mathbf{f}_{\text{ext}} = \boldsymbol{0}.
 \end{equation} 

\subsection{Linearization}
As in Section <REF>, \eqref{curvy_f_system} will be solved using Newton's method. This would require consistent linearisation of $\mathbf{f}_{\text{int}}$ and $\mathbf{f}_{\text{ext}}$ with respect to
 X AND WHAT??. The derivations of linearised quantities required to formulate $\Delta \mathbf{f}_{\text{int}}^{e}$ and $\Delta \mathbf{f}_{\text{ext}}^{e}$ can be found in the Appendix \ref{appendices:curvy_linearisations}.
\begin{equation}\label{eq:f_int}
\mathbf{f}_{\mathrm{int}}^{e} = \int_{\Omega^{e}} \mathbf{N}^{T}_{, \alpha} \tau^{\alpha \beta} \mathbf{N}_{, \beta} \, \mathrm{d} A \mathbf{x}_{e}
\end{equation}

\begin{equation}\label{eq:del_f_int}
    \begin{split}
    \Delta \mathbf{f}_{\mathrm{int}}^{e} &= \int_{\Omega^{e}} \mathbf{N}^{T}_{, \alpha} \Delta \tau^{\alpha \beta} \mathbf{N}_{, \beta} \, \mathrm{d} A \mathbf{x}_{e} + \int_{\Omega^{e}} \mathbf{N}_{, \alpha}^{T} \tau^{\alpha \beta} \mathbf{N}_{, \beta} \, \mathrm{d} A \Delta \mathbf{x}_{e} \\ &= \int_{\Omega^{e}} \mathbf{N}_{, \alpha}^{T} c^{\alpha \beta \gamma \delta} \boldsymbol{a}_{\gamma} \cdot \mathbf{N}_{\delta} \Delta \mathbf{x}_{e} \mathbf{N}_{, \beta} \, \mathrm{d} A \mathbf{x}_{e} + \int_{\Omega^{e}} \mathbf{N}_{, \alpha}^{T} \tau^{\alpha \beta} \mathbf{N}_{, \beta} \, \mathrm{d} A \Delta \mathbf{x}_{e} \\ &= \int_{\Omega^{e}} \mathbf{N}^{T}_{,\alpha} c^{\alpha \beta \gamma \delta} \left( \boldsymbol{a}_{\gamma} \cdot \mathbf{N}_{\delta} \right) \left( \mathbf{N}_{, \beta} \mathbf{x}_{e} \right) \, \mathrm{d} A \Delta \mathbf{x}_{e} + \int_{\Omega^{e}} \mathbf{N}_{, \alpha}^{T} \tau^{\alpha \beta} \mathbf{N}_{, \beta} \, \mathrm{d} A \Delta \mathbf{x}_{e} \\ &= \int_{\Omega^{e}} c^{\alpha \beta \gamma \delta} \mathbf{N}_{, \alpha}^{T} \left( \boldsymbol{a}_{\gamma} \mathbf{N}_{\delta} \right) \boldsymbol{a}_{\beta} \, \mathrm{d} A \Delta \mathbf{x}_{e} + \int_{\Omega^{e}} \mathbf{N}_{, alpha}^{T} \tau^{\alpha \beta} \mathbf{N}_{, \beta} \, \mathrm{d} A \Delta \mathbf{x}_{e} \\ &= \int_{\Omega^{e}} c^{\alpha \beta \gamma \delta} \mathbf{N}_{, \alpha}^{T} \left( \boldsymbol{a}_{\alpha} \otimes \boldsymbol{a}_{\gamma} \right) \mathbf{N}_{, \delta} \, \mathrm{d} A \Delta \mathbf{x}_{e} + \int_{\Omega^{e}} \mathbf{N}_{, \alpha}^{T} \tau^{\alpha \beta} \mathbf{N}_{, \beta} \, \mathrm{d} A \Delta \mathbf{x}_{e}
    \end{split}
\end{equation}

\begin{equation}\label{eq:del_f_int_2}
    \begin{split}
        \Delta \mathbf{f}_{\mathrm{int}}^{e} &= \int_{\Omega^{e}} \left[ c^{\alpha \beta \gamma \delta} \mathbf{N}_{, \alpha}^{T} \left( \boldsymbol{a}_{\beta} \otimes \boldsymbol{a}_{\gamma} \right) \mathbf{N}_{, \delta} + \mathbf{N}_{, \alpha}^{T} \tau^{\alpha \beta} \mathbf{N}_{, \beta} \right] \, \mathrm{d} A \Delta \mathbf{x}_{e} \\ &= \underbrace{\int_{\Omega^{e}} c^{\alpha \beta \gamma \delta} \mathbf{N}_{, \alpha}^{T} \left( \boldsymbol{a}_{\beta} \otimes \boldsymbol{a}_{\gamma} \right) \mathbf{N}_{, \delta} \, \mathrm{d} A}_{=: k_{\mathrm{mat}}^{e}} + \underbrace{\int_{\Omega^{e}} \mathbf{N}_{,\alpha}^{T} \tau^{\alpha \beta} \mathbf{N}_{, \beta} \, \mathrm{d} A}_{=: k_{\mathrm{geo}}^{e}}
    \end{split}
\end{equation}

\begin{equation}\label{eq:del_f_int_3}
\Delta \mathbf{f}_{\mathrm{int}}^{e} = \left( k_{\mathrm{mat}}^{e} + k_{\mathrm{geo}}^{e} \right) \Delta \mathbf{x}_{e}
\end{equation}

\begin{equation}\label{eq:f_ext}
    \mathbf{f}_{\mathrm{ext}}^{e} = \int_{\Omega^{e}} \mathbf{N}^{T} \boldsymbol{f}_{0} \, \mathrm{d} A + \int_{\partial_{t} \Omega^{e}} \mathbf{N}^{T} \overline{\boldsymbol{t}} \, \mathrm{d} s + \int_{\Omega^{e}} \mathbf{N}^{T} p \boldsymbol{n} \, \mathrm{d} a
\end{equation}

For dead loading~\eqref{eq:f_ext} can be written as

\begin{equation}\label{eq:f_ext_dead}
\mathbf{f}_{\mathrm{ext}}^{e} = \int_{\Omega^{e}} \mathbf{N}^{T} p \boldsymbol{n} \, \mathrm{d} a
\end{equation}

\begin{equation}\label{eq:del_f_ext_dead}
    \begin{split}
    \Delta \mathbf{f}_{\mathrm{ext}}^{e} &= \int_{\Omega^{e}} \mathbf{N}^{T} \Delta p \boldsymbol{n} \, \mathrm{d}a + \int_{\Omega^{e}} \mathbf{N}^{T} p \Delta \left( \boldsymbol{n}\, \mathrm{d} a \right) = \int_{\Omega^{e}} \mathbf{N}^{T} \Delta \left( \rho \boldsymbol{g} \cdot \boldsymbol{x} \right) \boldsymbol{n} \, \mathrm{d} a + \int_{\Omega^{e}} \mathbf{N}^{T} p \Delta \left( \boldsymbol{n} \, \mathrm{d} a \right) \\ &= \int_{\Omega^{e}} \mathbf{N}^{T} \rho \left( \underbrace{\Delta \boldsymbol{g} \cdot \boldsymbol{x}}_{= 0} + \boldsymbol{g} \cdot \Delta \boldsymbol{x}_{e} \right) \boldsymbol{n} \, \mathrm{d} a + \int_{\Omega^{e}} \mathbf{N}^{T} p \Delta \left( \boldsymbol{n} \, \mathrm{d} a \right) 
    \end{split}
\end{equation}

\begin{equation}\label{eq:del_f_ext_dead_2}
    \begin{split}
        \Delta \mathbf{f}_{\mathrm{ext}}^{e} &= \int_{\Omega^{e}} \rho \mathbf{N}^{T} \left( \boldsymbol{g} \cdot \Delta \mathbf{x}_{e} \right) \boldsymbol{n} \, \mathrm{d} a + \int_{\Omega^{e}} \mathbf{N}^{T} p \Delta \left( \boldsymbol{n} \, \mathrm{d} a \right) \\ &= \int_{\Omega^{e}} \rho \mathbf{N}^{T} \left( \boldsymbol{n} \otimes \boldsymbol{g} \right) \, \mathrm{d} a \Delta \mathbf{x}_{e} + \int_{\Omega^{e}} p \mathbf{N}^{T} \left( \boldsymbol{n} \otimes \boldsymbol{a}^{\alpha} - \boldsymbol{a}^{\alpha} \otimes \boldsymbol{n} \right) \mathbf{N}_{, \alpha} \, \mathrm{d} a \Delta \mathbf{x}_{e}
    \end{split}
\end{equation}

\end{document}